\section{Boundary phases and the specified suspension}\label{sec:phase}
There is another reason to keep separated boundary and bulk observables distinct. A relative spectral asymmetry can produce a boundary phase in one regularization. Its dependence on a change of local counterterms can be tested explicitly here. Write $g=|\Gamma|$, $P_3(x)=2x^3/3-x/2$, and use $[F]_*^R=F(R)-F(q_*)$. Below the first shell crossing the smooth part of the shifted intrinsic eta function is
\[
 \eta(D_0-sX)=\eta_0+sP_3(X)/g .
\]
One can obtain the polynomial from the identity-class spectrum on the cover: differentiate the continued eta series inside the gap and evaluate its residue, giving $\partial_X\eta=s(2X^2-1/2)/g$. Nonidentity elements of a free quotient have no local heat residue. Crossing a shell of multiplicity $m_s$ adds $-2s m_s$ on moving from the smaller to the larger $X$ side. At $h=0$, one-sided limits are therefore required. Subtracting charges $-\eta/2$ with a common local term at the throat recovers the integer in Theorem~\ref{thm:fermions}.

For the signed boundary-eta phase with the outer-minus-inner orientation, the branch difference modulo $2\pi$ is $-(\pi/g)[P_3(X)]_*^R$. At $h=0$, $R\to\infty$ this is $3\pi/(2g)$. Possible endpoint contacts have even complex multiplicity here and exponentiate to unity in this comparison. Let $B=sX/q$ be the signed normal axial component and $\mathcal R_3=6/q^2$. Each branch admits the same local boundary functional
\begin{equation}
 S_{\rm ct}(s)=
 \left[\int_{Y_q}\left(-\frac{B\mathcal R_3}{48\pi}
                   +\frac{B^3}{6\pi}\right)\dd V\right]_*^R
 =\frac{\pi s}{2g}[P_3(X)]_*^R .
 \label{eq:ct}
\end{equation}
The coefficients follow by inserting $\vol(Y_q)=2\pi^2q^3/g$. Subtracting the values on the two branches cancels the measured boundary phase at every $h,R,g$ in this model. One need not classify all allowed counterterms to exhibit this cancellation.

For a separate closed test geometry use a smooth positive round joining profile at the outer ends of the double. A direct mirror of the outgoing boundary has a kink if $q'\ne0$. Take $q\ge q_*$, strict on an open set, and keep $X=3c/(2q^2)$. On a product interval in a fifth coordinate let $f$ vary from $1$ to $-1$ within those bounds and replace $sX$ by $fX$. For each $\lambda$ the kernel multiplier over the four-dimensional time circle is $\exp[-\oint(\lambda-fX)\dd\tau/q]$, or its inverse. It is positive and different from unity, even on the extreme contacting modes. Periodic and antiperiodic time spin structures both exclude a four-dimensional kernel. The four-dimensional spectral flow is zero.

More is needed to deduce a phase from a suspension. With chirality $\Gamma_5=\mathrm{diag}(1,-1)$ use the product operator $\mathcal D_5=i\Gamma_5\partial_v+D_4(v)$ with the self-adjoint spectral domain at its two ends \cite{DF,WY}. In a shell basis the radial coefficients are real. Complex conjugation followed by $\Gamma_5$ anticommutes with $\mathcal D_5$ and commutes with the tangential boundary projectors of $-i\Gamma_5D_4$. Thus the APS spectrum is symmetric and $\eta_5(0)=0$. Each round three-dimensional eigenspace with a real spin lift has its usual quaternionic structure, giving even complex multiplicity also in the kernel of the suspension. Its exponentiated reduced eta phase $e^{-i\pi(\eta_5+\dim\ker\mathcal D_5)}$ is $+1$. This computation specifies the closing geometry, interpolation and boundary conditions; a general Weyl determinant line can carry additional data.
